\section{Summary}

This report describes the GPU port of the ALF auxiliary-field quantum Monte Carlo code that lives on the branch \texttt{gpu-port-unified}. It is built from the Markdown notes in \file{docs/gpu/} and from the verification runs of 2026-11-13. Every number comes from one of those notes. The source file is given in a comment next to the number in the \LaTeX{} source. Missing numbers and open items are marked in red as TODO.

\subsection*{What runs on which hardware}

One source tree builds for four targets (\file{README.md}, \file{BUILD.md}). The default build, \texttt{GPU=none}, leaves the CPU code unchanged.

\begin{table}[h]
\centering
\caption{Build targets and the hardware on which they were run.}
\label{tab:targets}
\begin{tabular}{@{}lp{4.6cm}p{3.6cm}p{4.2cm}@{}}
\toprule
\texttt{GPU=} & Hardware used & Compiler & Libraries \\
\midrule
\texttt{none} & any CPU & any ALF compiler & BLAS/LAPACK \\ % README.md
\texttt{nvidia} & A30, H100 & nvfortran 25.7 / 26.9 & cuBLAS, cuSOLVER \\ % README.md
\texttt{intel} & Max 1550 (SuperMUC-NG Phase 2) & ifx, oneAPI 2025.x & oneMKL, SYCL shim \\ % README.md
\texttt{amd} & MI210 & gfortran 14 (amdgcn offload) & hipBLAS, hipSOLVER, HIP kernels \\ % README.md
\bottomrule
\end{tabular}
\end{table}

The port was developed for the Hubbard model with a space-time point defect (time-dependent hopping, $N_{\rm FL}=2$). It runs the finite-temperature sweep, the local updates (standard, delayed and submatrix), the wraps, the stabilisation and the equal-time and time-displaced measurements on the device. Random numbers, the accept/reject decision, bookkeeping and file output stay on the host. MPI with independent chains and parallel parameter sets (\texttt{PARALLEL\_PARAMS}) work. Real parallel tempering, the projector code, Langevin, HMC and MALA, other Hamiltonians and several other features are refused at setup, and the run then uses the CPU path (Section~\ref{sec:ported}).

\subsection*{Headline speeds}

Table~\ref{tab:headline} gives seconds per sweep for the Hubbard defect model with $\beta=L$, $\Delta\tau=0.1$, submatrix update and \texttt{Nwrap} 10. Numbers are from different sessions and different timers. Section~\ref{sec:perf} lists how each was obtained.

\begin{table}[h]
\centering
\caption{Seconds per sweep, one chain, and aggregate throughput of the best measured layout at $L=24$.}
\label{tab:headline}
\begin{tabular}{@{}lrrl@{}}
\toprule
Platform and layout & $L=24$ & $L=32$ & Source \\
\midrule
A30, one chain & 5.25 & 18.2 & \file{TUNING\_DEFECT.md} \\ % TUNING_DEFECT.md
A30, 4 ranks with MPS, per chain & 10.5 & \todo{not measured} & \file{MPI\_GPU.md}, \file{README.md} \\ % README.md
H100, one chain & 3.51 & 9.82 & \file{helma\_h100.md} \\ % helma_h100.md 4.1, 4.2
H100, 8 ranks with MPS, per chain & 6.80 & 34.5 & \file{helma\_h100.md} \\ % helma_h100.md
Max 1550, one chain per tile (V7) & 4.95 & 16.4 & \file{ISSUE\_INTEL\_VS\_A30.md} \\ % ISSUE_INTEL_VS_A30.md sec 2
Max 1550, 4 chains per tile, per chain & 9.3 & \todo{not measured} & \file{ISSUE\_INTEL\_VS\_A30.md} \\ % ISSUE_INTEL_VS_A30.md, p7 NOTES
MI210, one process & 4.41 & 17.5 & \file{amd\_mi210/README.md} \\ % platforms/amd_mi210/README.md
\bottomrule
\end{tabular}
\end{table}

The main points:
\begin{itemize}
\item At $L=24$ all four GPUs are within a factor 1.5 of each other for one chain (3.51 to 5.25~s per sweep). The H100 is the fastest at both sizes. At $L=32$ the H100 is 1.8 times faster than the A30 and the MI210 and Max~1550 are close to the A30 (16.4 to 18.2~s).
\item One chain does not fill a GPU. Several chains per GPU raise the throughput by a factor 2 to 4 but slow every chain down by a factor 2 to 3. On the H100 at $L=24$ eight ranks deliver 1.18 sweeps per second against 0.285 for one rank (Section~\ref{sec:mps}).
\item Against one CPU thread the A30 is 28 times faster at $L=24$ and 9.6 times at $L=16$, and not faster at $L=12$ (1.76 against 1.68~s). A CPU comparison with a full node for the defect model does not exist \todo{CPU node-level number for the defect model}.
\item On Intel, run-time switches of one binary reduced the single-chain time at $L=24$ from 10.6 to 4.96~s per sweep (Section~\ref{sec:intel}). The two last steps, device QR and device \texttt{ZUNMQR}, are not yet released for production.
\item MPS is required whenever several ranks share an NVIDIA GPU. Without it, two ranks on an H100 are 6.5 times slower per chain at $L=24$ (Section~\ref{sec:mps}).
\end{itemize}

\subsection*{State of verification}

The GPU build is compared with a CPU build of the same compiler, same parameters and same seeds. Agreement is at rounding level, about $10^{-13}$ with \texttt{Mz = .F.} and about $10^{-10}$ with \texttt{Mz = .T.} (Section~\ref{sec:verif}). The consolidation of the six vendor branches into one was checked on 2026-11-13:
\begin{itemize}
\item \texttt{GPU=none}: four fixed-seed cases are bitwise identical to the pre-merge tip \texttt{cc473abc}. All tested upstream features are bitwise identical to upstream ALF master \texttt{521ddfca} when \texttt{LOBS\_ST = 1} is set. Two differences from upstream are intended.
\item AMD MI210: bitwise identical to the pre-merge AMD tip \texttt{156a0b01} at $L=6$, 12, 16 and 24; timing unchanged.
\item Intel Max~1550 (ifx 2025.3, JIT and AOT builds): against the same-compiler CPU build the deviation is $2.2\times10^{-15}$ and $1.5\times10^{-14}$ ($L=6$ and 12, \texttt{Mz = .F.}) and $5.5\times10^{-11}$ ($L=6$, \texttt{Mz = .T.}). Against the pre-merge Intel build, Green functions, scalars, acceptance and configurations are bitwise identical. With the V7 switches the $L=6$, \texttt{Mz = .T.} value is $1.36\times10^{-10}$, the same as before the merge, above the tolerance of $10^{-10}$.
\item NVIDIA A30: the code builds with nvfortran 26.9; the GPU runs are \todo{pending} (card occupied). NVIDIA H100 (Helma): \todo{pending}.
\end{itemize}

\subsection*{Open points that affect production}

The Intel steps V4 and V7 wait for a numerics review of the device QR. The $L=6$, \texttt{Mz = .T.} test fails its pre-registered tolerance with V7 ($1.36\times10^{-10}$ against $10^{-10}$). The race test of the device QR cannot detect a race on the \texttt{tau} slot. The staging-buffer setting, worth 33\,\% at $L=24$ on Intel, is not in the production job script of the running set. Section~\ref{sec:limits} lists these and the decisions that belong to the owner.
